% part: intuitionistic-logic
% chapter: semantics
% section: propositions

\documentclass[../../../include/open-logic-section]{subfiles}

\begin{document}

\olfileid{int}{sem}{pro}

\olsection{विधानसंच}

कधीकधी !!a{relational model}~$\mModel{M}$ मधील ज्या
जगतांत !!a{formula}~$!A$ सत्य असते, त्या जगांच्या संचाला नाव देणे
उपयुक्त ठरते. त्याला $!A$ ने परिभाषित केलेला
\emph{विधानसंच} म्हणू आणि $\Prop{M}{!A}$ असे लिहू.

\begin{defn}
  $V$ या फलनाचा विस्तार $\Prop{M}{!A} \colon \Frm
  \to \Pow{W}$ अशा फलनात पुढीलप्रमाणे विगमनाने करू:
  \begin{enumerate}
  \item $\Prop{M}{\lfalse} = \emptyset$
  \item $\Prop{M}{p} = V(p)$
  \item $\Prop{M}{\lnot !B} = {}$
    \[
      \Setabs{w \in W}{\text{$Rww'$ असेल अशा प्रत्येक $w'$ साठी $w'
          \notin \Prop{M}{!B}$}}.
    \]
  \item $\Prop{M}{!B \land !C} = \Prop{M}{!B} \cap \Prop{M}{!C}$.
  \item $\Prop{M}{!B \lor !C} = \Prop{M}{!B} \cup \Prop{M}{!C}$.
  \item $\Prop{M}{!B \lif !C} = {}$
    \[
      \Setabs{w \in W}{\text{$Rww'$ असेल अशा प्रत्येक $w'$ साठी,
          $w' \in \Prop{M}{!B}$ असेल तर $w' \in \Prop{M}{!C}$}}.
    \]
  \end{enumerate}
\end{defn}
%READERNOTE{OLINC-331: स्रोतामध्ये Prop(M,A) या एका मूल्याला संपूर्ण सूत्रांवरील फलनाचे नाव दिले आहे; येथे स्रोताचे चिन्ह जतन केले असून आशय A पासून जगांचा संच देणारा नकाशा असा वाचावा.}

\begin{prop}
  \ollabel{prop:propositions-true}
  $\mSat{M}{!A}[w]$ तेव्हा आणि केवळ तेव्हाच $w \in \Prop{M}{!A}$.
\end{prop}
%READERNOTE{OLINC-330: स्रोताने पूर्ति-संबंधाच्या जागी फक्त प्रतिमानाचे नाव छापणारा mModel मॅक्रो वापरला आहे; येथे mSat केला आहे.}

\begin{proof}
  सराव.
\end{proof}

\begin{prob}
  \olref[int][sem][pro]{prop:propositions-true} सिद्ध करा.
\end{prob}

\end{document}
